\documentclass{article}

% packages
\usepackage{amsmath, amsthm, thmtools, amsfonts, amssymb, luacode, catchfile, tikzducks, hyperref, ifthen}
\ifcsname c@kobocompile\endcsname
	\usepackage[a5paper, total={1072pt, 1448pt}, margin=10pt, includeheadfoot]{geometry} % set page margins
\else
	\usepackage[a4paper, margin=50pt, includeheadfoot]{geometry}
\fi
\usepackage[shortlabels]{enumitem}
\usepackage[skip=3pt, indent=0pt]{parskip}

% language
\usepackage[bidi=basic, layout=tabular, provide=*]{babel}
\ifcsname c@english\endcsname
	\babelprovide[main, import]{english}
\else
	\babelprovide[main, import]{hebrew}
	\babelprovide{rl}
\fi
%\babelfont{rm}{Libertinus Serif}
\babelfont{rm}[Renderer=Harfbuzz]{Libertinus Serif}
\babelfont{sf}{Libertinus Sans}
\babelfont{tt}{Libertinus Mono}

% style
\AddToHook{cmd/section/before}{\clearpage}	% Add line break before section
\linespread{1.3}
\setcounter{secnumdepth}{0}		% Remove default number tags from sections, this won't do well with theorems
\AtBeginDocument{\setlength{\belowdisplayskip}{3pt}}
\AtBeginDocument{\setlength{\abovedisplayskip}{3pt}}
\graphicspath{ {../images/} }

% operators
\DeclareMathOperator\cis{cis}
\DeclareMathOperator\Sp{Sp}
\DeclareMathOperator\tr{tr}
\DeclareMathOperator\im{Im}
\DeclareMathOperator\re{Re}
\DeclareMathOperator\diag{diag}
\DeclareMathOperator*\lowlim{\underline{lim}}
\DeclareMathOperator*\uplim{\overline{lim}}
\DeclareMathOperator\rng{rng}
\DeclareMathOperator\Sym{Sym}
\DeclareMathOperator\Arg{Arg}
\DeclareMathOperator\Log{Log}
\DeclareMathOperator\dom{dom}
\DeclareMathOperator\supp{Supp}
\DeclareMathOperator\var{Var}
\DeclareMathOperator\cov{Cov}

% commands
%\renewcommand\qedsymbol{\textbf{מש''ל}}
%\renewcommand\qedsymbol{\fbox{\emoji{lizard}}}
\newcommand{\Aa}[0]{\mathcal{A}}
\newcommand{\Bb}[0]{\mathcal{B}}
\newcommand{\CC}[0]{\mathbb{C}}
\newcommand{\Cc}[0]{\mathcal{C}}
\newcommand{\EE}[0]{\mathbb{E}}
\newcommand{\FF}[0]{\mathbb{F}}
\newcommand{\Ff}[0]{\mathcal{F}}
\newcommand{\Ii}[0]{\mathcal{I}}
\newcommand{\Gg}[0]{\mathcal{G}}
\newcommand{\Ll}[0]{\mathcal{L}}
\newcommand{\Mm}[0]{\mathcal{M}}
\newcommand{\NN}[0]{\mathbb{N}}
\newcommand{\Nn}[0]{\mathcal{N}}
\newcommand{\PP}[0]{\mathbb{P}}
\newcommand{\Pp}[0]{\mathcal{P}}
\newcommand{\QQ}[0]{\mathbb{Q}}
\newcommand{\RR}[0]{\mathbb{R}}
\newcommand{\Rr}[0]{\mathcal{R}}
\newcommand{\Ss}[0]{\mathcal{S}}
\newcommand{\TT}[0]{\mathbb{T}}
\newcommand{\Uu}[0]{\mathcal{U}}
\newcommand{\Vv}[0]{\mathcal{V}}
\newcommand{\Ww}[0]{\mathcal{W}}
\newcommand{\ZZ}[0]{\mathbb{Z}}
\newcommand{\acts}[0]{\circlearrowright}
\newcommand{\explain}[2] {
	\begin{flalign*}
		 && \text{#2} && \text{#1}
	\end{flalign*}
}
\newcommand{\maketitleprint}[0]{ \begin{center}
	%\begin{tikzpicture}[scale=3]
	%	\duck[graduate=gray!20!black, tassel=red!70!black]
	%\end{tikzpicture}	
	\includegraphics[width=6cm]{cover}
\end{center}
}

% theorem commands
\newtheoremstyle{c_remark}
	{}	% Space above
	{}	% Space below
	{}% Body font
	{}	% Indent amount
	{\bfseries}	% Theorem head font
	{}	% Punctuation after theorem head
	{.5em}	% Space after theorem head
	{\thmname{#1}\thmnumber{ #2}\thmnote{ \normalfont{\text{(#3)}}}}	% head content
\newtheoremstyle{c_definition}
	{3pt}	% Space above
	{3pt}	% Space below
	{}% Body font
	{}	% Indent amount
	{\bfseries}	% Theorem head font
	{}	% Punctuation after theorem head
	{.5em}	% Space after theorem head
	{\thmname{#1}\thmnumber{ #2}\thmnote{ \normalfont{\text{(#3)}}}}	% head content
\newtheoremstyle{c_plain}
	{3pt}	% Space above
	{3pt}	% Space below
	{\itshape}% Body font
	{}	% Indent amount
	{\bfseries}	% Theorem head font
	{}	% Punctuation after theorem head
	{.5em}	% Space after theorem head
	{\thmname{#1}\thmnumber{ #2}\thmnote{ \text{(#3)}}}	% head content

\ifcsname c@english\endcsname
	\theoremstyle{plain}
	\newtheorem{theorem}{Theorem}[section]
	\newtheorem{lemma}[theorem]{Lemma}
	\newtheorem{proposition}[theorem]{Proposition}
	\newtheorem*{proposition*}{Proposition}
	%\newtheorem{corollary}[theorem]{אין חלופה עברית}

	\theoremstyle{definition}
	\newtheorem{definition}[theorem]{Definition}
	\newtheorem*{definition*}{Definition}
	\newtheorem{example}{Example}[section]
	\newtheorem{exercise}{Exercise}[section]

	\theoremstyle{remark}
	\newtheorem*{remark}{Remark}
	\newtheorem*{solution}{Solution}
	\newtheorem{conclusion}[theorem]{Conclusion}
	\newtheorem{notation}[theorem]{Notation}
\else
	\theoremstyle{c_plain}
	\newtheorem{theorem}{משפט}[section]
	\newtheorem{lemma}[theorem]{למה}
	\newtheorem{proposition}[theorem]{טענה}
	\newtheorem*{proposition*}{טענה}
	%\newtheorem{corollary}[theorem]{אין חלופה עברית}

	\theoremstyle{c_definition}
	\newtheorem{definition}[theorem]{הגדרה}
	\newtheorem*{definition*}{הגדרה}
	\newtheorem{example}{דוגמה}[section]
	\newtheorem{exercise}{תרגיל}[section]

	\theoremstyle{c_remark}
	\newtheorem*{remark}{הערה}
	\newtheorem*{solution}{פתרון}
	\newtheorem{conclusion}[theorem]{מסקנה}
	\newtheorem{notation}[theorem]{סימון}
\fi

% Questions related commands
\newcounter{question}
\setcounter{question}{1}
\newcounter{sub_question}
\setcounter{sub_question}{1}

\ifcsname c@english\endcsname
	\newcommand{\question}[1][0]{
		\ifthenelse{#1 = 0}{}{\setcounter{question}{#1}}
		\section{Question \arabic{question}}
		\addtocounter{question}{1}
		\setcounter{sub_question}{1}
	}

	\newcommand{\subquestion}[1][0]{
		\ifthenelse{#1 = 0}{}{\setcounter{sub_question}{#1}}
		\subsection{Part \alph{sub_question}}
		\addtocounter{sub_question}{1}
	}
\else
	\newcommand{\question}[1][0]{
		\ifthenelse{#1 = 0}{}{\setcounter{question}{#1}}
		\section{שאלה \arabic{question}}
		\addtocounter{question}{1}
		\setcounter{sub_question}{1}
	}

	\newcommand{\subquestion}[1][0]{
		\ifthenelse{#1 = 0}{}{\setcounter{sub_question}{#1}}
		\subsection{סעיף \localecounter{letters.gershayim}{sub_question}}
		\addtocounter{sub_question}{1}
	}
\fi

% import lua and start of document
\directlua{common = require ('../common')}

\GetEnv{AUTHOR}

% headers
\author{\AUTHOR}
\date\today

\title{פתרון מטלה 03 --- חשבון אינפינטסמלי 3 (80415)}

\begin{document}
\maketitle
\maketitleprint{}

\Question{}
נמצא שתי קבוצות סגורות ולא ריקות $A, B \subseteq \RR^2$ כך שהמרחק ביניהן לא מתקבל. \\*
נבחר את הקבוצות $A = \{ (x, y) \in \RR^2 \mid y = 0\}, B = \{ (x, y) \in \RR^2 \mid x y = 1 \}$. \\*
שתי הקבוצות ניתנות לתיאור על־ידי פונקציות $\RR \to \RR$ רציפות בכל תחומן ולכן הקבוצות מכילות את כל הנקודות הגבוליות שלהן, ובהתאם סגורות. \\*
לעומת זאת, הקבוצות זרות, ונראה כי לכל $x \in \RR$ מתקיים $\rho((x, 0), (x, \frac{1}{x})) = \sqrt{0^2 + \frac{1}{x^2}} = \frac{1}{|x|}$ ולכן $\inf \text{dist}(A, B) = 0$

\Question{}
יהי $V$ מרחב וקטורי ו־$\lVert \cdot \rVert, \lVert \cdot \rVert'$ שתי נורמות עליו. \\*
תהי $S_1 = \{ v \in V \mid \lVert v \rVert = 1 \}$ ונגדיר $A = \{ \lVert v\rVert' \mid v \in S_1\}$. \\*
נוכיח כי הנורמות שקולות אם ורק אם $\inf A > 0, \sup A < \infty$.
\begin{proof}
	\textbf{כיוון ראשון:}
	נניח כי הנורמות שקולות. \\*
	לכן קיימים $0 < C_1 \le C_2$ כך ש־$\forall v \in V : C_1 \lVert v \rVert \le \lVert v \rVert' \le C_2 \lVert v \rVert$. \\*
	בפרט
	\[
		\forall v \in S_1 : C_1 \le \lVert v \rVert' \le C_2
		\implies \forall a \in A : C_1 \le a \le C_2
	\]
	ובהתאם $\inf A \ge C_1 \ge 0$ וגם $\sup A \le C_2 < \infty$.

	\textbf{כיוון שני:}
	נניח כי $\inf A = a > 0$ ו־$\sup A = b < \infty$. \\*
	לכן נובע ישירות $\forall x \in A : a \le x \le b$ ובהתאם $\forall v \in S_1 : a \lVert v \rVert \le \lVert v \rVert' \le b \lVert v \rVert$. \\*
	יהי $u \in V$, אז אנו יודעים כי קיים $u^* \in V$ כך ש־$\lVert u^* \rVert = 1$ ו־$\lambda \in \RR^+$ כך ש־$u = \lambda u^*$ ולכן גם $\lVert u \rVert = \lambda \lVert u^* \rVert = \lambda$. \\*
	נשים לב עתה כי $u^* \in S_1$ ולכן $a \lVert u^* \rVert \le \lVert u^* \rVert' \le b \lVert u^* \rVert$, נכפיל את שלושת האגפים ב־$\lambda$ ונקבל
	\[
		a \lVert u \rVert \le \lVert u \rVert' \le b \lVert u \rVert
	\]
	ולכן הנורמות שקולות.
\end{proof}

\Question{}
יהי $(V, \lVert \cdot \rVert)$ מרחב נורמי ממימד $d$ מעל $\RR$.

\Subquestion{}
תהי $T : \RR^d \to V$ העתקה לינארית חד־חד ערכית ועל. \\*
נוכיח שהפונקציה $f : \RR^d \to \RR$ המוגדרת על־ידי $f(u) = \lVert T(u) \rVert$ משרה נורמה על $\RR^d$.
\begin{proof}
	נראה כי שלוש התכונות של נורמה מתקיימות.
	\begin{enumerate}
		\item חיוביות: ידוע כי $T$ הפיכה ולכן גרעינה הוא וקטור האפס בלבד, ובהתאם $f(u) = \lVert Tu \rVert = 0 \iff \lVert u \rVert = 0$.
		\item כפל בסקלר: $\forall u \in \RR^2, \lambda \in \RR, f(\lambda u) = \lVert T (\lambda u) \rVert  = \lVert \lambda \cdot Tu \rVert = |\lambda| \lVert Tu \rVert = |\lambda| f(u)$.
		\item אי־שוויון המשולש: $\forall u, v \in \RR^2, f(u + v) = \lVert T(u + v) \rVert = \lVert Tu + Tv \rVert \le \lVert T u \rVert + \lVert T v \rVert = f(u) + f(v)$.
	\end{enumerate}
	ומצאנו כי $f$ מקיימת את התנאים לנורמה ולכן מהווה נורמה ל־$\RR^d$.
\end{proof}
יהי $B = (v_1, \dots, v_d)$ בסיס ל־$V$, ונגדיר נורמה על־ידי
\[
	\lVert \alpha v_1 + \cdots + \alpha_d v_d \rVert_1 = \sum_{i = 1}^{d} |\alpha_1| = d |\alpha_1|
\]
נוכיח שהנורמות $\lVert \cdot \rVert, \lVert \cdot \rVert_1$ שקולות.
\begin{proof}
	בסעיף הקודם הוכחנו ששימוש ב־$T$ מאפשר להשרות נורמה מ־$V$ ל־$\RR^2$, ולכן אפשר להגדיר פונצקיה משרה לנורמה לשתי הנורמות הנתונות. \\*
	נשתמש בטענה כי כל הנורמות ב־$\RR^d$ הן שקולות ונקבל כי גם $\lVert \cdot \rVert$ שקולה ל־$\lVert \cdot \rVert_1$.
\end{proof}

\Question{}
יהיו $X, Y$ מרחביים מטריים ו־$f : X \to Y$ פונקציה חד־חד ערכית ועל. \\*
נוכיח כי שלושת התנאים הבאים שקולים:
\begin{enumerate}
	\item $f$ היא הומיאומורפיזם.
	\item כל קבוצה $U \subseteq X$ היא פתוחה אם ורק אם $f(U)$ פתוחה.
	\item כל קבוצה $U \subseteq X$ היא סגורה אם ורק אם $f(U)$ סגורה.
\end{enumerate}

\begin{proof}
	\textbf{1 \leftarrow{} 2:}
	נניח כי $f$ היא הומיאומורפיזם. \\*
	לכן $f$ רציפה וגם $f^{-1}$ רציפה. לכן נסיק מהתנאי המספיק לרציפות כי $U \subseteq X$ פתוחה אז $f(U)$ פתוחה ולכל $U \subseteq Y$ גם $f^{-1}(U) \subseteq X$ פתוחה. \\*
	במסתכם $U \subseteq X$ פתוחה אם ורק אם $f(U) \subseteq Y$ פתוחה.

	\textbf{2 \leftarrow{} 3:}
	נניח כי $U \subseteq X$ פתוחה אם ורק אם $f(U) \subseteq Y$ פתוחה. \\*
	תהי קבוצה סגורה $A \subseteq X$, אז $A^C$ קבוצה פתוחה ולכן $f(A^C)$ קבוצה פתוחה ב־$Y$. \\*
	באופן דומה $f(U)$ קבוצה פתוחה. אילו נניח בשלילה כי $f(U) \cap f(U^C) \ne \emptyset$ נקבל כי ישנם $x \ne y \in X$ כך ש־$f(x) = f(y)$ בסתירה לחד־חד ערכיות, ולכן $f(U) \cap f(U^C) = \emptyset$. \\*
	נראה גם $f(U) \cup f(U^C) = Y$ אחרת תכונת העל נסתרת עבור $f$, ולכן נסיק $f(U^C) = {(f(U))}^C$. \\*
	מטענה זו נוכל להסיק ישירות ש־$U \subseteq X$ סגורה אם ורק אם $f(U) \subseteq Y$ סגורה.

	\textbf{2 \leftarrow{} 3:}
	נניח כי כל קבוצה $U \subseteq X$ סגורה אם ורק אם $f(U) \subseteq Y$ סגורה אף היא. \\*
	נוכל להשתמש בחד־חד ערכיות באופן זהה לסעיף הקודם וכך גם בעל ונקבל כי $f(U^C) = {(f(U))}^C$ לכל $U \subseteq X$. \\*
	נשים לב כי התנאי המספיק לרציפות, שמקור כל קבוצה פתוחה הוא קבוצה פתוחה מתקיים לשני הכיוונים, לכן $f, f^{-1}$ רציפות ובהתאם $f$ הומיאומורפיזם.
\end{proof}

\Question{}
יהיו $X, Y$ מרחבים מטריים ו־$f : X \to Y$ פונקציה חד־חד ערכית, על ורציפה.

\Subquestion{}
נוכיח שאם $X$ קומפקטי אז $f$ היא הומיאומורפיזם.
\begin{proof}
	ידוע ש־$X$ קומפקטית ו־$f$ רציפה ולכן נסיק ש־$f(X) \subseteq Y$ קומפקטית אף היא. \\*
	ידוע ש־$f$ על ולכן $f(X) = Y$ ו־$Y$ קבוצה קומפקטית. \\*
	עוד ידוע שתת־קבוצה סגורה של קבוצה קומפקטית היא קומפקטית, יהי $A \subseteq X$ סגור, אז גם $f(A) \subseteq Y$ קבוצה סגורה. \\*
	נניח ש־$f(B)$ קבוצה סגורה כאשר $B \subseteq X$. לכן גם $f(B)$ היא קומפקטית, אם נניח ש־$B$ עצמה לא קומפקטית נקבל סתירה לטענה שהוצגה קודם, לכן $B$ קומפקטית. \\*
	$B \subseteq X$ וקומפקטית ולכן גם סגורה, וקיבלנו כי תנאי 3 של השאלה הקודמת מתקיים ונובע כי $f$ הומיאומורפיזם.
\end{proof}

\Subquestion{}
נראה דוגמה למקרה בו $X$ לא קומפקטית ו־$f$ איננה הומיאומורפיזם. \\*
נבחן את הפונקציה מאחד התרגילים הקודמים $f : [0, 1) \to S(0, 1) \subseteq \RR^2$ המוגדרת על־ידי % chktex 9
\[
	f(t) = (\cos t, \sin t)
\]
אבל הראינו כי היא לא הומיאומורפיזם בשאלה 6 סעיף ה'.

\Question{}
\Subquestion{}
יהי $X$ מרחב מטרי שלם, ויהי $D \subseteq X$ כדור סגור. \\*
נוכיח ש־$D$ עם המטריקה המושרית על־ידי $X$ היא מרחב מטרי שלם.
\begin{proof}
	תהי סדרת קושי ${(a_n)}_{n = 1}^\infty \subseteq D$. \\*
	כמובן שגם $\forall n \in \NN a_n \in X$ ולכן מהשלמות נסיק ש־$(a)$ מתכנסת ונגדיר $\lim_{n \to \infty} a_n = a$. \\*
	ידוע כי $a \in A$, אבל ידוע גם כי  $D$ קבוצה סגורה ולכן היא מכילה את כלל הנקודות הגבוליות שלה, ואכן $a$ נקודה גבולית שלה ונסיק כי $a \in D$. \\*
	מכאן נסיק כי $D$ עצמה שלמה.
\end{proof}

\Subquestion{}
נוכיח שלמשוואה
\[
	f'(x) = x \cdot \cos (f^2(x))
\]
עם תנאי ההתחלה $f(0) = 3$ יש פתרון יחיד בסביבת $0$.
\begin{proof}
	תחילה נמיר למשוואה אינטגרלית ונקבל
	\[
		f(x) = 3 + \int_{0}^{x} t \cdot \cos(f^2(t)) dt
	\]
	נבחר קטע $I = [-a, a]$ סביב $0$, ונגדיר $T : C(I) \to C(I)$ על־ידי
	\[
		T(f)(x) = 3 + \int_{0}^{x} t \cdot \cos(f^2(t)) dt
	\]
	יהיו $f, g \in C(I)$ ונניח $\lVert f - g \rVert_\infty = M$, ויהי $x \in I$, מתקיים
	\begin{align*}
		| (Tf)(x) - (Tg)(x) |
		& = | \int_{0}^{x} t \cdot \cos(f^2(t)) dt - \int_{0}^{x} t \cdot \cos(g^2(t)) dt | \\
		& = | \int_{0}^{x} t \cdot (\cos(f^2(t)) - \cos(g^2(t))) dt | \\
		& \le | \int_{0}^{x} t M dt | \\
		& \le M | \frac{x^2}{2} | \\
		| (Tf)(x) - (Tg)(x) |
		& \le \frac{a^2 M}{2} \\
	\end{align*}
	ונסיק כי ההעתקה $T$ מכווצת עבור כל $0 < a < \sqrt{2}$ שנבחר. \\*
	תנאי משפט נקודת השבת של בנך מתקיימים ונקבל כי קיימת נקודה $x_0 \in I$ יחידה המקיימת $(Tf)(x_0) = x_0$ ועל־פי הגדרת ההעתקה זהו גם הפתרון של המשוואה הנתונה.
\end{proof}

\Question{}
תהי $I \subseteq \RR$ סביבה של $0$ ותהי $f : I \to \RR$ פונקציה גזירה כך ש־$f(0) = 0$ ו־$f'$ רציפה ב־$0$. נניח גם ש־$|f'(0)| < 1$. \\*
נוכיח שקיים קטע $X = [-a, a]$ עבורו $f(X) \subseteq X$ וגם ההעתקה המצומצמת $f : X \to X$ היא מכווצת ממש.
\begin{proof}
	מהנתון $|f'(0)| < 1$ ו־$f'$ רציפה בסביבת $0$ נסיק כי קיימת סביבה $X = [-a, a]$ כך ש־$\forall x \in X : |f'(x)| < 1$. \\*
	נגדיר $M = \max_{x \in X} |f'(x)|$. כמובן $M < 1$ על־פי הגדרתו. \\*
	נקבל כי $f^*(x) = |f(x)| - x$ מקיימת $f^*(x) \le 0$ עבור $0 \le x \le a$ וממעבר דומה על החלק השלילי של התחום שבחרנו לקבל ש־$f(X) \subseteq X$.

	נגדיר $g : X \to X$ על־ידי $g(x) = f(x)$ ההעתקה המצומצמת. \\*
	ראינו כי $|x| < a \implies |g(x)| < |x|$, דהינו ההעתקה מכווצת, ואף $|g(x)| < M|x|$ בתחום.
	\[
		\forall x, y \in (-a, a) : |g(x) - g(y)| \le |g(x)| - |g(y)| < M|x| - M|y|
		\forall x, y \in (-a, a) : |g(x) - g(y)| \le \left| M|x| - M|y| \right|
	\]
	ותמשיך מפה נמאס לי
\end{proof}

\end{document} % chktex 17
